\documentclass[11pt,a4paper]{article}

\usepackage[a4paper,margin=1in]{geometry}
\usepackage{fontspec}
\IfFontExistsTF{Times New Roman}{\setmainfont{Times New Roman}}{\setmainfont{TeX Gyre Termes}}
\IfFontExistsTF{Menlo}{\setmonofont{Menlo}}{\setmonofont{Latin Modern Mono}}
\usepackage{amsmath,amssymb}
\usepackage{booktabs}
\usepackage{longtable}
\usepackage{array}
\usepackage{enumitem}
\usepackage{graphicx}
\usepackage{fancyhdr}
\usepackage[round,authoryear]{natbib}
\usepackage[hidelinks]{hyperref}
\usepackage{xurl}
\usepackage{microtype}

\setlength{\parindent}{0pt}
\setlength{\parskip}{0.65em}
\setlist[itemize]{leftmargin=1.4em}
\newcommand{\file}[1]{\path{#1}}
\newcommand{\code}[1]{\path{#1}}
\newcolumntype{L}[1]{>{\raggedright\arraybackslash}p{#1}}

\pagestyle{fancy}
\fancyhf{}
\lhead{\small Pilot 01 Report}
\rhead{\small Controlled Stress-Map Diagnostics}
\cfoot{\thepage}
\renewcommand{\headrulewidth}{0.4pt}

\title{\textbf{Pilot 01 Report: Controlled Pilot Stress Maps}\\
\large Validated Baseline American CN/PSOR Solver}
\author{Codex-assisted pilot diagnostics report}
\date{June 17, 2026}

\begin{document}
\maketitle

\begin{abstract}
This report documents the first controlled pilot stress-map experiment after the solver
validation ladder and downstream design gate. The pilot uses the baseline American
Crank--Nicolson / projected SOR solver only. It produces small diagnostic tables and figures for
an American put family and a dividend-paying American call family. The outputs are pilot
diagnostics, not production risk surfaces, training labels, broad parameter sweeps, or neural
surrogate inputs.
\end{abstract}

\tableofcontents
\newpage

\section{Purpose of the Pilot}

Ticket 12 concluded \code{PASS_SOLVER_VALIDATION_WITH_LIMITATIONS}. The downstream design gate
then approved a small, human-reviewed pilot before any large-scale stress maps or data generation.
This Pilot 01 report records that first narrow experiment.

The purpose is practical: verify that the validated classical solver can produce interpretable
stress-map style figures together with the metadata needed to read them safely. The pilot keeps
the scope deliberately small. It checks option values, payoffs, continuation premium,
exercise/continuation indicators, approximate boundary curves, PSOR convergence, LCP diagnostics,
and masked Greek diagnostics.

\section{Link to Prior Gates}

The pilot follows the validation path documented in the Ticket 12 synthesis report and the
downstream design gate. The theory basis remains the classical Black--Scholes PDE and American
finite-difference/LCP setting \citep{black_scholes_1973,merton_1973,brennan_schwartz_1977,
wilmott_howison_dewynne_1995}. The Rannacher comparison from Ticket 10A remains a comparison
tool only \citep{rannacher_1984}; it is not used as the pilot default.

\begin{longtable}{L{0.30\textwidth} L{0.54\textwidth}}
\toprule
Gate decision & Pilot implementation rule \\
\midrule
\endhead
Use validated baseline solver & All runs use \code{american_crank_nicolson_psor_price}. \\
Keep Rannacher optional & No Pilot 01 run uses the Rannacher solver. \\
Use normalized strike & Every run uses \(K=1\). \\
Use practical pilot grid & Main runs use \(S_{\max}=4\), \(M=N=120\). \\
Include higher-grid check & Base put and base dividend call also run at \(M=N=180\). \\
Keep outputs diagnostic & Manifest status is \code{pilot_diagnostic_only}. \\
\bottomrule
\end{longtable}

\section{Solver and Grid Choice}

The default solver is the Ticket 04 baseline American CN/PSOR solver. The pilot uses
\[
S_{\max}=4K,\qquad M=N=120
\]
for the main cases and \(M=N=180\) for the two higher-grid checks. The main interpretation region
is \(S/K \in [0.4,1.8]\). The solver still computes on the full domain because the boundary
conditions and diagnostics depend on the full grid.

\section{Pilot Case Definitions}

\begin{longtable}{L{0.25\textwidth} L{0.22\textwidth} L{0.38\textwidth}}
\toprule
Family & Base case & One-at-a-time variations \\
\midrule
\endhead
American put & \(T=1\), \(r=0.05\), \(q=0.02\), \(\sigma=0.20\), \(S_{\max}=4\), \(M=N=120\). & \(\sigma=0.40\), \(r=0.03\), \(q=0.00\), plus \(M=N=180\) base check. \\
Dividend-paying American call & \(T=1\), \(r=0.05\), \(q=0.08\), \(\sigma=0.20\), \(S_{\max}=4\), \(M=N=120\). & \(\sigma=0.40\), \(r=0.03\), \(q=0.03\), \(q=0.10\), plus \(M=N=180\) base check. \\
\bottomrule
\end{longtable}

This is intentionally not a broad parameter sweep. It is a one-at-a-time pilot that a human
reviewer can inspect.

\section{Output Definitions}

\begin{longtable}{L{0.28\textwidth} L{0.52\textwidth}}
\toprule
Output & Meaning \\
\midrule
\endhead
Value grid \(U(S,\tau)\) & American option value from baseline CN/PSOR. \\
Payoff \(\Phi(S)\) & Immediate exercise payoff. \\
Continuation premium \(U-\Phi\) & Diagnostic amount by which the value exceeds payoff. \\
Exercise indicator & \(1\) where \(U-\Phi \le 10^{-6}\), otherwise \(0\). \\
Boundary curve & Threshold-based extraction from continuation premium. \\
LCP diagnostics & Obstacle, equation, and complementarity checks from Ticket 08 logic. \\
Delta/Gamma diagnostics & Finite-difference Greeks with boundary, kink, maturity, and strict masks. \\
\bottomrule
\end{longtable}

\section{Diagnostic Requirements}

Every run records solver, grid, parameter, boundary, and diagnostic metadata. The acceptance rule
for pilot review is conservative:
\[
\text{PSOR converged},\quad
\max(\Phi-U,0)\le 10^{-8},\quad
\max(b-AU,0)\le 10^{-6},\quad
\max |(U-\Phi)(AU-b)|\le 10^{-6}.
\]

All 11 Pilot 01 runs met this rule and were marked
\code{ACCEPTABLE_FOR_PILOT_REVIEW}. This does not make the figures production risk surfaces. It
means the small pilot is numerically coherent enough for human review.

\section{Tables and Interpretation}

\begin{longtable}{L{0.30\textwidth} L{0.18\textwidth} L{0.18\textwidth} L{0.20\textwidth}}
\toprule
Observed item & American put base & Dividend call base & Interpretation \\
\midrule
\endhead
PSOR convergence & True & True & All time steps converged. \\
Max obstacle violation & \(0\) & \(0\) & Payoff obstacle respected. \\
Max equation violation & \(1.55664789392\times10^{-9}\) & \(2.14416621448\times10^{-9}\) & Near diagnostic tolerance. \\
Max complementarity product & \(6.6483476629\times10^{-11}\) & \(2.37960951462\times10^{-11}\) & Small LCP product residual. \\
Boundary found rows & 120 & 120 & Boundaries found for all positive-\(\tau\) rows. \\
Strict max \(|\Gamma|\) & \(3.18001886582\) & \(2.6798873723\) & Greek diagnostics remain cautionary. \\
\bottomrule
\end{longtable}

\begin{longtable}{L{0.26\textwidth} L{0.22\textwidth} L{0.22\textwidth} L{0.20\textwidth}}
\toprule
Higher-grid check & \(M=N=120\) boundary range & \(M=N=180\) boundary range & Reading \\
\midrule
\endhead
American put base & \(0.766754\) to \(0.934082\) & \(0.777814\) to \(0.955653\) & Boundary estimates shift with grid and remain approximate. \\
Dividend call base & \(1.066165\) to \(1.266630\) & \(1.044356\) to \(1.266615\) & Higher grid is a check, not a final convergence proof. \\
\bottomrule
\end{longtable}

The generated tables are:
\begin{itemize}
  \item \file{results/02_pilot_stress_maps/tables/pilot_01_run_summary.csv};
  \item \file{results/02_pilot_stress_maps/tables/pilot_01_diagnostic_summary.csv};
  \item \file{results/02_pilot_stress_maps/tables/pilot_01_boundary_summary.csv};
  \item \file{results/02_pilot_stress_maps/tables/pilot_01_selected_slices.csv};
  \item \file{results/02_pilot_stress_maps/tables/pilot_01_output_manifest.csv}.
\end{itemize}

\section{Figures and Interpretation}

Figure~\ref{fig:value_heatmaps} shows option values over the main interpretation region. These
are the most direct pilot stress maps, but they should be read together with diagnostics rather
than as standalone production surfaces.

\begin{figure}[htbp]
\centering
\includegraphics[width=0.95\textwidth]{results/02_pilot_stress_maps/figures/pilot_01_base_value_heatmaps.png}
\caption{Base American put and dividend-paying American call value heatmaps.}
\label{fig:value_heatmaps}
\end{figure}

Figure~\ref{fig:premium_heatmaps} shows continuation premium \(U-\Phi\). The premium is central
to obstacle interpretation and to the boundary extraction, but it is not the same thing as
American-minus-European value.

\begin{figure}[htbp]
\centering
\includegraphics[width=0.95\textwidth]{results/02_pilot_stress_maps/figures/pilot_01_base_premium_heatmaps.png}
\caption{Base continuation-premium heatmaps.}
\label{fig:premium_heatmaps}
\end{figure}

Figure~\ref{fig:indicator_maps} shows the exercise/continuation indicator with boundary overlay.
The map is threshold-based: it visualizes the diagnostic rule \(U-\Phi \le 10^{-6}\), not an exact
analytical exercise region.

\begin{figure}[htbp]
\centering
\includegraphics[width=0.95\textwidth]{results/02_pilot_stress_maps/figures/pilot_01_base_indicator_boundary_maps.png}
\caption{Exercise/continuation indicator maps with approximate boundary overlays.}
\label{fig:indicator_maps}
\end{figure}

Figure~\ref{fig:value_slices} gives selected value slices. These are easier for a greenhand reader
to inspect than heatmaps because each curve can be compared with payoff.

\begin{figure}[htbp]
\centering
\includegraphics[width=0.95\textwidth]{results/02_pilot_stress_maps/figures/pilot_01_base_selected_value_slices.png}
\caption{Selected value slices at \(\tau/T=0.01,0.25,0.50,0.75,1.00\).}
\label{fig:value_slices}
\end{figure}

Figure~\ref{fig:boundary_variations} compares boundary curves across one-at-a-time variations.
This is the most useful pilot figure for seeing how dividend yield, interest rate, and volatility
move the threshold boundary. It remains a diagnostic boundary curve, not an exact free-boundary
claim.

\begin{figure}[htbp]
\centering
\includegraphics[width=0.95\textwidth]{results/02_pilot_stress_maps/figures/pilot_01_boundary_variation_comparison.png}
\caption{Approximate boundary curves across base and one-at-a-time variations.}
\label{fig:boundary_variations}
\end{figure}

Figure~\ref{fig:premium_variations} shows mid-time continuation-premium slices across variations.
It helps connect the boundary curves to the underlying premium quantity.

\begin{figure}[htbp]
\centering
\includegraphics[width=0.95\textwidth]{results/02_pilot_stress_maps/figures/pilot_01_premium_slice_variation_comparison.png}
\caption{Continuation-premium slices across variations near \(\tau/T=0.5\).}
\label{fig:premium_variations}
\end{figure}

Figure~\ref{fig:greek_slices} shows Gamma diagnostics for the two base cases. The plotted Gamma
curves include caution markers for payoff kinks and extracted boundary context. They should not be
used as production Greek labels.

\begin{figure}[htbp]
\centering
\includegraphics[width=0.95\textwidth]{results/02_pilot_stress_maps/figures/pilot_01_greek_diagnostic_slices.png}
\caption{Base Gamma diagnostic slices with kink and boundary cautions.}
\label{fig:greek_slices}
\end{figure}

\section{Acceptance Criteria Results}

\begin{longtable}{L{0.42\textwidth} L{0.18\textwidth} L{0.28\textwidth}}
\toprule
Criterion & Result & Interpretation \\
\midrule
\endhead
All PSOR steps converge or report non-convergence & PASS & All 11 runs converged. \\
Obstacle violation near zero & PASS & All runs report maximum obstacle violation \(0\). \\
LCP diagnostics near tolerance & PASS & Largest equation violation is \(1.93971868349\times10^{-8}\); complementarity remains small. \\
Boundary status recorded & PASS & Every run records found/not-found counts and boundary status. \\
Greek masks recorded & PASS & Selected slices and diagnostic summaries include masks and strict metrics. \\
Full parameter metadata present & PASS & Every output table records case parameters and grid metadata. \\
No downstream data products created & PASS & Outputs are CSV summaries and figures only. \\
\bottomrule
\end{longtable}

\section{Limitations}

\begin{itemize}
  \item The pilot uses a small one-at-a-time case set. It is not a broad stress-map sweep.
  \item Boundary extraction uses a fixed \(10^{-6}\) continuation-premium threshold.
  \item The \(M=N=180\) checks are confirmation checks, not a formal convergence proof.
  \item Gamma remains fragile near payoff kinks, maturity rows, and boundary bands.
  \item Figures are pilot diagnostics. They are not production risk surfaces.
  \item No output should be treated as a training target for a neural surrogate.
\end{itemize}

\section{What This Pilot Supports}

The pilot supports a narrow claim: the validated baseline solver can generate small,
metadata-rich diagnostic stress-map figures for representative American put and dividend-paying
American call cases. It also supports using the same table/manifest structure for a future
human-reviewed pilot expansion.

\section{What This Pilot Does Not Support}

The pilot does not support large-scale data generation, production Greek reporting, neural
surrogate training, exact analytical free-boundary claims, or final paper-style claims about broad
parameter regimes. Those steps require a separate approved plan and further validation.

\section{Recommended Next Step}

The recommended next step is human review of the Pilot 01 report, figures, and manifest. If the
review accepts the scope and caveats, the next implementation should still be modest: a slightly
expanded pilot with the same metadata contract, not a broad dataset or neural surrogate stage.

\section{Reproducibility Record}

\begin{longtable}{L{0.32\textwidth} L{0.52\textwidth}}
\toprule
Item & Record \\
\midrule
\endhead
Source module & \file{src/american_risk_surfaces/downstream/pilot_stress_maps.py}. \\
Experiment script & \file{experiments/11_pilot_stress_maps.py}. \\
Test file & \file{tests/test_pilot_stress_maps.py}. \\
Report source & \file{reports/04_downstream/pilot_stress_maps_report.tex}. \\
Report PDF & \file{reports/04_downstream/pilot_stress_maps_report.pdf}. \\
Output folder & \file{results/02_pilot_stress_maps/}. \\
Bibliography & \file{reports/03_solver/references.bib}. \\
\bottomrule
\end{longtable}

\bibliographystyle{plainnat}
\bibliography{reports/03_solver/references}

\end{document}
